\chapter{Organización del club ambiental}

El club ambiental busca ser un lugar donde todos los integrantes puedan aportar a las actividades de manera equitativa, por lo que se buscará mantener una organización democrática, tomando las decisiones con la cooperación de todos los integrantes.

Todos los integrantes del club podrán proponer actividades, y tanto la mejora como la organización de estas podrán ser llevadas a cabo por todos los integrantes del club, así como también el participar en tales actividades. El proponer actividades y participar en ellas también es posible para estudiantes que no participen de manera constante en el club.

Los materiales de trabajo, háblese de manuales, encuestas, datos, etcétera, se espera que se mantengan abiertos a todo estudiantes que guste consultarlos, esto con el objetivo de que puedan conocer las actividades del club y su funcionamiento incluso si no tienen la disponibilidad de asistir constantemente a este. La edición sin embargo de los materiales de trabajo estará reservada para los estudiantes participantes del club.

\section{Mantenimiento del centro de composta}

El mantenimiento de un centro de compostaje necesita de personas que, además de estar capacitadas en el tema, se dediquen activamente a ello, por ende, se buscará definir un grupo de alumnos voluntarios que se dediquen principalmente al mantenimiento de la composta, los cuales se buscará que tengan una disponibilidad de horario adecuada para las necesidades de la composta. La participación de los integrantes del club en el mantenimiento de la composta no estará limitado aún así a este grupo designado.

Habría dos actividades principales a tratar en el centro de compostaje:rRecolección de residuos y el mantenimiento del estado de la composta.

\subsection{Recolección de basura}

Los encargados de esta área serán responsables de ir periódicamente a los contenedores de basura orgánica para recoger su contenido, transportarlo al centro de compostaje y separarlos residuos para que estén listos para su deposición.

\subsection{Mantenimiento de la composta}

Los responsables de esta área se encargarán de triturar y depositar la basura recolectada en el compostador que se encuentre activo en el momento. Tanto antes como después de esto tomarán las mediciones diarias correspondientes (pH, temperatura, etc.) y las subirán a la base de datos del club.
